% =============================================================================
%  Scholarly-maturity pages
%
%    Lay Summary
%    Statement of Originality
%    Statement of Contributions  (CRediT taxonomy)
%    Reproducibility & Data-Availability Statement
% =============================================================================

% -----------------------------------------------------------------------------
%  Lay Summary  (plain-language abstract for non-experts)
% -----------------------------------------------------------------------------
\cleardoublepage
\chapter*{Lay Summary}
\addcontentsline{toc}{chapter}{Lay Summary}
\thispagestyle{plain}
\markboth{Lay Summary}{}

\begingroup
\setlength{\parindent}{0pt}

\dropcap{R}{ailway} tracks rest on a layer of crushed stone called
\emph{ballast}. The ballast keeps the rails in place, spreads the wheel
loads to the soil beneath, and lets rainwater drain away. Every time a
train passes, the ballast is pushed and bumped, and the stones slowly
break down. Once the ballast becomes too crushed, the track sinks, the
ride becomes rough, and the railway must spend money and time on
maintenance.

\bigskip

This thesis asks a simple question: \emph{how quickly does ballast break
down under the conditions seen on Indian Railways, and what can we do
to slow it down?} The answer matters because the Indian Railways
network is the largest passenger railway in the world; even a small
improvement in ballast life saves substantial maintenance cost and
keeps trains running on time.

\bigskip

To answer the question, the authors built a heavy laboratory machine
that drops a 50-kilogram hammer onto a tub full of ballast,
reproducing the kind of impact that a train wheel gives to a rough
section of track. By placing different rubber mats and different
cushions underneath the ballast, the machine can show which
combinations protect the stones best. By scanning forty-five
individual ballast stones in three dimensions both \emph{before} and
\emph{after} they are pounded, the authors can measure---for the
first time on Indian ballast---exactly how the shape of each stone
changes.

\bigskip

The thesis reports two main discoveries. First, polyurethane rubber
mats placed under the ballast reduce the amount of breakage
significantly, especially when the base of the track is rigid (as it
is over bridges and culverts). Second, ballast stones break in
\emph{two distinct ways}: most stones lose their sharp corners
gradually (a slow, manageable process), but a small fraction
\emph{shatter} catastrophically---and the shattering happens only on
rigid bases. This second finding suggests that the
hardest-to-maintain stretches of Indian Railways are the bridges and
the transition zones at their ends, and that under-ballast mats
should be deployed there as a priority.

\endgroup

% -----------------------------------------------------------------------------
%  Statement of Originality
% -----------------------------------------------------------------------------
\cleardoublepage
\chapter*{Statement of Originality}
\addcontentsline{toc}{chapter}{Statement of Originality}
\thispagestyle{plain}
\markboth{Statement of Originality}{}

\begingroup
\setlength{\parindent}{0pt}

We, \textbf{Durgesh Kumar Legha} (Roll No.\ 2270007) and
\textbf{Sandeep Patel} (Roll No.\ 2280043), declare that this thesis
titled \emph{Assessment of Railroad Ballast Degradation under Modified
Drop-Weight Impact Loading Conditions and Image-Based Morphological
Analysis} is our own original work, conducted between July~2024 and
April~2026 at the School of Civil Engineering, Gati Shakti
Vishwavidyalaya, Vadodara, under the supervision of
Dr.~Dinesh~Gundavaram.

\bigskip

We further declare that:

\begin{enumerate}[leftmargin=*,itemsep=4pt]
    \item All published and unpublished sources used in this thesis are
          cited in the text and listed in the References. The
          BibTeX-based bibliography contains every reference cited;
          there are no uncited references and no unreferenced
          citations.
    \item All laboratory data, including the nine raw sieve-analysis
          spreadsheets, the in-line load-cell records, and the
          ninety paired \gls{STL} meshes (pre- and post-test, for the
          forty-five tracer particles), were generated by the authors
          using the apparatus described in \cref{chap:methodology}.
          No data have been fabricated, falsified, or selectively
          omitted.
    \item The novel contributions of this thesis (the modified
          drop-weight apparatus, the per-particle paired protocol, the
          Mode-A / Mode-B damage taxonomy, and the
          \texttt{ballast\_analyzer.py} pipeline) are, to the best of
          our knowledge, original to this work. Closely related prior
          work is critically reviewed in \cref{chap:literature}.
    \item No part of this thesis has been submitted previously, either
          in whole or in part, for any degree at any university.
    \item The two journal manuscripts under review and the two
          accepted Springer conference papers listed in
          \emph{Outcomes of This Project} are derived from the work
          reported here; the relevant chapters cite those manuscripts
          where overlap occurs.
\end{enumerate}

\vspace{1.5cm}
\noindent
\begin{tabular}{p{0.42\textwidth}p{0.42\textwidth}}
\rule{0pt}{40pt}\hrulefill & \rule{0pt}{40pt}\hrulefill \\
\textbf{Durgesh Kumar Legha} & \textbf{Sandeep Patel} \\
Roll No.\ 2270007            & Roll No.\ 2280043 \\
\end{tabular}

\vspace{1.0cm}
\noindent Vadodara, April 2026.

\endgroup

% -----------------------------------------------------------------------------
%  Statement of Contributions (CRediT taxonomy)
% -----------------------------------------------------------------------------
\cleardoublepage
\chapter*{Statement of Contributions}
\addcontentsline{toc}{chapter}{Statement of Contributions}
\thispagestyle{plain}
\markboth{Statement of Contributions}{}

\begingroup
\setlength{\parindent}{0pt}

The work reported in this thesis was jointly carried out by the two
candidates under the supervision of Dr.~Dinesh~Gundavaram. The
following contribution matrix follows the
\href{https://credit.niso.org/}{CRediT (Contributor Roles Taxonomy)} of
the National Information Standards Organization, which is
increasingly adopted by international journals to attribute author
contributions transparently.

\bigskip

{\footnotesize
\renewcommand{\arraystretch}{1.18}
\rowcolors{2}{rowbg}{white}
\begin{tabularx}{\textwidth}{|>{\raggedright\arraybackslash}p{0.32\textwidth}|>{\centering\arraybackslash}X|>{\centering\arraybackslash}X|>{\centering\arraybackslash}X|}
\hline
\headrow
\textbf{\color{white}CRediT role} & \textbf{\color{white}D.~K.~Legha} &
\textbf{\color{white}S.~Patel} & \textbf{\color{white}D.~Gundavaram} \\
\hline
Conceptualisation             & supporting & supporting & lead       \\
Methodology                   & lead       & lead       & supporting \\
Apparatus design \& fabrication & lead     & equal      & supporting \\
Software (\texttt{ballast\_analyzer.py}) & equal & lead & supporting \\
Validation                    & equal      & equal      & supporting \\
Formal analysis               & lead       & equal      & supporting \\
Investigation (experiments)   & equal      & equal      & supporting \\
Resources / funding           & ---        & ---        & lead       \\
Data curation                 & equal      & lead       & supporting \\
Writing -- original draft     & lead       & equal      & supporting \\
Writing -- review \& editing  & equal      & equal      & lead       \\
Visualisation (TikZ, plots)   & lead       & equal      & supporting \\
Supervision                   & ---        & ---        & lead       \\
Project administration        & supporting & supporting & lead       \\
\hline
\end{tabularx}}

\vspace{14pt}

\noindent
Where two candidates are marked \emph{lead} or \emph{equal} for the same
role, the work was performed jointly and is not separable into
individual products. Where one candidate is marked \emph{lead}, that
candidate took primary responsibility for the role and is the principal
defender of that aspect of the work.

\vspace{14pt}

\noindent
\textbf{\color{gsvblue}Chapter-by-chapter contributions.}
\Cref{chap:introduction,chap:literature} were drafted by D.~K.~Legha
and revised jointly. The apparatus design and instrumentation
documented in \cref{chap:methodology} are due principally to
D.~K.~Legha (mechanical design, drop tower) and S.~Patel (load-cell
integration, instrumentation, structured-light scanning protocol).
\Cref{chap:impact-results} (bulk results) was drafted by D.~K.~Legha;
\cref{chap:morph-results} (per-particle morphology) was drafted by
S.~Patel. \Cref{chap:conclusions} was drafted jointly. The
\texttt{ballast\_analyzer.py} pipeline reproduced in
\cref{chap:appB} was written principally by S.~Patel with code review
and unit-test design by D.~K.~Legha.

\endgroup

% -----------------------------------------------------------------------------
%  Reproducibility & Data-Availability Statement
% -----------------------------------------------------------------------------
\cleardoublepage
\chapter*{Reproducibility \& Data-Availability Statement}
\addcontentsline{toc}{chapter}{Reproducibility \& Data Availability}
\thispagestyle{plain}
\markboth{Reproducibility \& Data Availability}{}

\begingroup
\setlength{\parindent}{0pt}

This thesis adheres to the reproducibility principles of the
\textbf{FAIR Guiding Principles} (Findable, Accessible, Interoperable,
Reusable) and reports its methods in sufficient detail that an
independent laboratory should be able to reproduce both the apparatus
and the analysis pipeline.

\bigskip

\begin{description}[leftmargin=2.6cm,style=nextline,
                    font=\bfseries\color{gsvblue}]

  \item[Apparatus] All structural drawings, instrumentation diagrams,
    and calibration certificates for the modified drop-weight impact
    apparatus are archived at the School of Civil Engineering, Gati
    Shakti Vishwavidyalaya, Vadodara. The apparatus is the subject of
    an Indian utility-patent application; the structural geometry
    and the calibration procedure are reproduced in
    \cref{chap:methodology} in sufficient detail that the apparatus
    can be replicated.

  \item[Materials] The basalt ballast is sourced from a single
    quarry near Vadodara; the gradation conforms to the Indian
    Railways \gls{RDSO} specification for \SI{65}{\milli\metre}-nominal
    ballast. The two under-ballast mats (\gls{PUBM} and \gls{RUBM})
    are commercially available and identified by manufacturer in
    \cref{sec:meth-PUBM,sec:meth-RUBM}. The flexible-base sand cushion
    and rigid-base steel plate are off-the-shelf laboratory items.

  \item[Raw data] The complete laboratory record comprises:
    \begin{enumerate}
       \item nine pre-test sieve-analysis spreadsheets and nine
             post-test sieve-analysis spreadsheets;
       \item nine in-line load-cell time-series files (one per test);
       \item ninety paired \gls{STL} meshes (\(45 \times 2\)) from
             the EinScan SL structured-light scanner;
       \item the master experimental log (Excel) cross-referencing
             the systematic test labels \testlabel{T1}--\testlabel{T9}
             to the legacy laboratory labels (\acrshort{LH},
             \acrshort{HH}).
    \end{enumerate}
    All raw data are archived on the GSV Civil Engineering server
    and on a redundant offline backup. The data are available from
    the authors on reasonable request for non-commercial research
    use, in accordance with the GSV data-sharing policy.

  \item[Software] The morphological-analysis pipeline
    (\texttt{ballast\_analyzer.py}, Python~3.10+, MIT licence) is
    reproduced in full in \cref{chap:appB} and depends on the
    publicly available \texttt{trimesh}, \texttt{numpy}, and
    \texttt{scipy} libraries. The pipeline has been unit-tested
    against synthetic reference meshes (sphere, ellipsoid, cube) for
    which the descriptors have known closed-form values; the test
    suite is included in the script docstring.

  \item[Analysis pipeline] All statistical tests reported in
    \cref{chap:morph-results} (Wilcoxon signed-rank test,
    Shapiro--Wilk normality test, Tukey $1.5\times$\gls{IQR}
    outlier-detection criterion) are reproducible from the raw
    paired meshes using the pipeline. Random seeds are not used;
    the analysis is deterministic.

  \item[Compilation] The thesis itself is compiled from the LaTeX
    source on the GSV Civil Engineering server with the command
    chain \texttt{pdflatex \(\to\) bibtex \(\to\) makeglossaries
    \(\to\) pdflatex \(\to\) pdflatex} (or equivalently
    \texttt{latexmk -pdf -bibtex}). The compile is deterministic
    when {\ttfamily SOURCE\_DATE\_EPOCH} is set.

\end{description}

\endgroup
